%! Fundamental Theorem of Algebra and Complex Zeros — Unit 4, Lesson 4.5 — Student Packet Cover Sheet
\documentclass[10pt]{article}
\usepackage{algebra2-article}
\usepackage{algebra2-boxes}
\usepackage{ltablex}
\keepXColumns

\begin{document}

% Full-bleed forest banner
\begin{tikzpicture}[remember picture, overlay]
  \fill[forest] (current page.north west)
    rectangle ([yshift=-0.9in]current page.north east);
\end{tikzpicture}
\vspace{-0.8in}

\begin{center}
  {\color{white}\textbf{\LARGE Algebra 2: Shepherd}}\\[4pt]
  {\color{forestmid}\large\textbf{Unit 4 \quad Polynomial Functions}}\\[6pt]
  {\color{white}\textbf{\large Lesson 4.5 \quad Fundamental Theorem of Algebra \& Complex Zeros}}
\end{center}

\vspace{0.22in}
\namedateperiod
\vspace{0.18in}

\begin{learningtargetbox}
\small
\begin{enumerate}[leftmargin=1.5em, itemsep=5pt]
  \item \ldots{} state the \textbf{Fundamental Theorem of Algebra} and use its counting rule: a polynomial of degree $n$ has \textbf{exactly $n$} complex zeros, counted with \textbf{multiplicity}.
  \item \ldots{} determine the \textbf{number and type} of solutions --- how many real, how many imaginary --- from the degree and the $x$-intercepts of a graph, \emph{before} solving.
  \item \ldots{} explain why imaginary zeros of a polynomial with \emph{real} coefficients always come in \textbf{complex conjugate pairs} $a+bi$ and $a-bi$.
  \item \ldots{} solve a polynomial equation of degree $3$ or higher \textbf{over the complex numbers}, finishing the depressed quadratic even when its discriminant is negative.
  \item \ldots{} write a polynomial in \textbf{factored form and standard form} given its zeros, supplying any missing conjugate.
  \item \ldots{} \textbf{verify} my solution list by counting it against the degree and by explaining which zeros a graph can and cannot show.
\end{enumerate}
\end{learningtargetbox}

\vspace{0.14in}

\begin{tocbox}
\small
\renewcommand{\arraystretch}{1.5}
\begin{tabularx}{\linewidth}{@{} c l X r @{}}
\rowcolor{forestbg}
\textbf{\#} & \textbf{Component} & \textbf{Description} & \textbf{Score} \\
\midrule
1 & Warm-Up        & Conjugate products, $x^2+4=0$, and finishing yesterday's cubic          & \blank{1.2cm} \\
2 & Guided Notes   & The counting theorem, conjugate pairs, solving over $\mathbb{C}$, building & \blank{1.2cm} \\
3 & Group Activity & \emph{Nothing Is Missing} --- counting, solving, and building, by tier   & \blank{1.2cm} \\
4 & Exit Ticket    & Quick check: one solve, one build, one SOL-style item                   & \blank{1.2cm} \\
5 & Homework       & Counting, three solves over $\mathbb{C}$, and polynomials to order       & \blank{1.2cm} \\
\midrule
  & \multicolumn{2}{r}{\textbf{Total}} & \blank{1.2cm} \\
\end{tabularx}
\end{tocbox}

\vspace{0.10in}

\begin{remindbox}
\small All year, ``no solution'' has really meant \textbf{no \emph{real} solution}: $x^2=-4$ in
Lesson 3.3, a negative discriminant in 3.5, a depressed polynomial that stalled yesterday. Today that
sentence gets finished. The \textbf{Fundamental Theorem of Algebra} says a polynomial of degree $n$
has \emph{exactly} $n$ complex zeros once you count repeated ones as many times as they repeat --- so
the degree tells you how many answers to look for before you find any of them. And when the
coefficients are real, imaginary zeros are never alone: if $a+bi$ is a zero, so is $a-bi$. That is why
imaginary zeros always come in an \emph{even} number, why an odd-degree graph must cross the $x$-axis,
and why a graph showing fewer $x$-intercepts than its degree is not hiding a mistake --- it is hiding
a conjugate pair, which no graph can show.
\end{remindbox}

\end{document}
